%% kapa: exam=1 round=37 year=2026 period=2 form=A subject=accounting subject_ko=회계학 seq=11 number=51
%% source: Q-Net 감정평가사 기출문제 (cst00309) artlSeq=5258924 file=raw/1008/5258924/2026년 제37회 감정평가사 1차 2교시.pdf
%% conversion: kordoc 4.24.51 -> tools/gen1.py (auto); review=proofread-sonnet; flags=none
\begin{question}{51}
(주)감평은 20×1년 초에 사채(액면금액 ￦100,000, 표시이자율 연 7 \%, 이자는 매년 말 지급, 만기 3년)를 발행하려고 하였으나 3개월이 경과된 20×1년 4월 1일에 동 사채를 발행하였다. 20×1년 초와 발행일 현재 유효이자율은 연 10 \%로 동일할 때, (주)감평이 20×1년 인식할 이자비용은? (단, 이자는 월할계산하고, 화폐금액은 소수점 첫째자리에서 반올림하며, 단수차이로 인한 오차가 있으면 가장 근사치를 선택한다.)
\par
{\small\centering
\begin{tabular}{|c|c|c|c|c|}\hline
\multirow{2}{*}{기간} & \multicolumn{2}{c|}{단일금액 ￦1의 현재가치} & \multicolumn{2}{c|}{정상연금 ￦1의 현재가치} \\\cline{2-5}
 & 7 \% & 10 \% & 7 \% & 10 \% \\\hline
3 & 0.8163 & 0.7513 & 2.6243 & 2.4868 \\\hline
\end{tabular}\par}
\begin{choices}
  \item ￦2,254
  \item ￦5,250
  \item ￦6,940
  \item ￦9,254
  \item ￦10,000
\end{choices}
\end{question}
\answer{3}
